\section{The every-occurrence spin bound}\label{ten:spin}
We prove Theorem~\ref{s:spin-moment}. The proof first isolates the sparse representations and then inducts on a balanced split of the coordinates. The signed tensor and stabilizer comparisons of Section~\ref{sec:static} supply the two local inputs to the induction. Put $J_n(\lambda)=D_\lambda\operatorname{Tr}_{V_\lambda}(B_n^*B_n)^r$, with the fixed integer $r$ still to be chosen.

\subsection{Dimensions outside the sparse range}
The induction needs a lower bound on dimension to absorb its counting
errors. Shapes with more than $n/2$ rows are already annihilated by
Lemma~\ref{lem:annihilation}; the following estimate treats the others.

\begin{lemma}\label{lem:spin-dimension-range}
There is an absolute $c>0$ such that, for all sufficiently large $n$,
if $\lambda\vdash n$ has at most $n/2$ rows and $k=n-\lambda_1\ge1$, then
\[
 \log D_\lambda\ge
 \begin{cases}
 c k\log(n/k),&k\le n/4,\\
 c n,&k\ge n/4.
 \end{cases}
\]
\end{lemma}
\begin{proof}
For $k\le n/4$, \eqref{eq:near-row-hooks} gives
$D_\lambda\ge e^{-1}\binom nk$, and the binomial lower bound proves
the assertion. Suppose $k\ge n/4$. If the first row and first column
both have length less than $n/10$, every hook is at most $n/5$, so the
hook formula gives $D_\lambda\ge n!/(n/5)^n$, exponential in $n$.
Otherwise, transpose if necessary to obtain a first row of length
$L\in[n/10,3n/4]$. Retain this row and $j=\lfloor n/100\rfloor$
boxes below it. These boxes can be chosen as a subdiagram, and its
tableaux extend to $\lambda$. Applying \eqref{eq:near-row-hooks} to
this subdiagram gives a lower bound
$e^{-1}\binom{L+j}{j}$, again exponential in $n$.
\end{proof}

\subsection{The sparse weighted-forest estimate}
Here is the operator estimate for small levels. For any fixed \(0<\delta<1\), for \(1\le k=n-\lambda_1\le n^{1-\delta}\),
\begin{equation}
 \|B_n\|_{V_\lambda}\le e^{O(k)} (2dk/n)^{k/4}
 \label{ten:spin:eq:5}
\end{equation}
for large \(n\), where the notation on the left means operator norm in \(V_\lambda\).

Let $\Omega_k$ be the ordered $k$-tuples of distinct positions. By
branching, $V_\lambda$ occurs in $\ell^2(\Omega_k)$ and is orthogonal
to every function that omits at least one tuple coordinate. We may
use either the sweep kernel or its adjoint, since their norms agree.
For $A\subseteq\{1,\ldots,k\}$ and any tuple $x$ whose
$A$-coordinates are distinct, define
\[
 F_A(x,y)=n^{|A|}\PP\{B_n(x_i)=y_i\text{ for all }i\in A\},
 \qquad L_A(x,y)=n^{-k}F_A(x,y).
\]
In this probability $B_n$ denotes a random sweep permutation.
Allowing arbitrary $y\in(\mathbb F_2^d)^k$, $L_A(x,\cdot)$ is the
probability law obtained by running all labels in $A$ through one
sweep and sending each other label independently to a uniform
position. Thus $L_A$ is a Markov kernel on tuples whose $A$-coordinates
are distinct; its restriction to $\Omega_k$ can lose mass.
In particular, $L_{\{1,\ldots,k\}}$ is the tuple kernel.

Given endpoints $x,y$, each label has a unique path through a sweep,
because each coordinate is processed once. Join two labels when
their paths use the same switch location. If label $i$ is isolated,
its switch prescriptions are independent of those of all other
labels. Hence $F_{A\cup\{i\}}(x,y)=F_A(x,y)$ whenever $i\notin A$.
On $V_\lambda$ we can therefore compute matrix coefficients with
\[
 n^{-k}\sum_A (-1)^{k-|A|} F_A(x,y)
\]
since all proper subsets give zero there. It vanishes if there is an isolated vertex. Consequently we bound the norm by the sum of norms of nonnegative
kernels
\[
 K_A(x,y)=n^{-k}F_A(x,y)
                 {\bf1}_{\{\text{the encounter graph has no isolated vertex}\}}.
\]

We next bound $K_A$ by a two-trip experiment. Draw $y$ from
$L_A(x,\cdot)$. From $y$, run a fresh reversed sweep for the labels
in $A$, and an independent fresh reversed sweep for each other label;
call the resulting tuple $x'$. For distinct $y,x'$, the transition
weight of this second trip is $L_A(x',y)$: the coupled tuple sweep
is replaced by its adjoint, and the one-label weights remain $1/n$.
Let $\mathcal E_2$ be the event that the second-trip paths have no
isolated label. Removing the first-trip encounter requirement and
both distinctness requirements gives
\[
 \sum_{x'\in\Omega_k}(K_AK_A^*)(x,x')
 =\sum_{y,x'\in\Omega_k}K_A(x,y)K_A(x',y)
 \le\PP_x(\mathcal E_2).
\]
The coupled labels remain distinct in either trip; the dropped
constraints concern the full tuple, including the independent labels.

Condition only on the switch settings of the second trip. There is
one fixed permutation network for $A$ and one for each label outside
$A$; the first-trip endpoints $y$ remain random with their original
law. Make a graph with one vertex for every starting position in
each of these networks. Join distinct vertices when their fixed
paths use the same switch location at a layer, also across networks.
A vertex has at most $2d$ neighbors in any one network.

Let $S$ be the set of the $k$ occupied vertices. For every set $T$
of distinct graph vertices we have
\[
 \PP\{T\subseteq S\}\le\prod_{w\in T}p_w,
 \qquad
 p_w=\begin{cases}|A|/n,&w\text{ in the }A\text{-network},\\
                  1/n,&w\text{ in a one-label network}.
       \end{cases}
\]
To see the assertion for the $A$-network, track its position subset
during the first trip. Initially it is fixed, so the joint-inclusion
bound by products of the one-site marginals holds. A fair independent
swap preserves this property. A test set containing exactly one
swapped site uses linearity of the two probabilities; for a set
containing both, the old product $p_up_v$ is at most the new product
$((p_u+p_v)/2)^2$. Sets containing neither are unchanged.
At the end every marginal is $|A|/n$, because each label is uniform.
The occupied starting positions in the one-label networks are
independent of this subset and of one another; selecting two vertices
in one such network has probability zero. This proves the displayed
joint-inclusion bound.
For a deterministic starting subset, this inclusion bound is also a
consequence of the stronger negative-dependence theory for random
swap-or-retain operations~\cite[Theorems~4.20 and~4.9]{BorceaBrandenLiggett2009}.

In this weighted graph the total mass is \(k\), and the weighted neighbor sum is bounded by \(\Delta=2dk/n\). To have no isolated vertices, the \(k\) occupied vertices must contain a forest on \(k\) vertices with component sizes at least two. Counting rooted plane trees for the components, the sum of product weights for candidates is bounded by
\[
 e^{O(k)}\sum_{1\le j\le k/2}\frac{k^j}{j!}\Delta^{k-j}.
\]
Indeed the plane forest shapes (ordered components) are at most exponentially many, and for any ordered shape the weight sum is bounded by \(k^j\Delta^{k-j}\) by descending from its \(j\) roots. One can divide by \(j!\) for the possible orders, since witnessing vertices are distinct. In the sparse range $\Delta\le1$ for large $n$, and
\[
 \sum_{1\le j\le k/2}\frac{k^j}{j!}\Delta^{k-j}
 \le\Delta^{k/2}\sum_{j\ge0}\frac{k^j}{j!}
 =e^k\Delta^{k/2}.
\]
For $k=1$ the encounter event is empty. For the symmetric nonnegative matrix $K_AK_A^*$, the maximum row
sum bounds the spectral radius. Therefore
\[
 \|K_A\|_{\mathrm{op}}^2
 \le\|K_AK_A^*\|_{\infty\to\infty}
 \le e^{Ck}\Delta^{k/2}.
\]
Taking square roots and summing the $2^k$ subset kernels gives
$e^{O(k)}\Delta^{k/4}$, proving \eqref{ten:spin:eq:5}.

\subsection{Parameters and the induction range}
We now turn to \eqref{ten:spin:eq:3}. Take a fixed small \(\epsilon\), say \(1/64\), and use
\[
 \eta_d=\bar\eta(1-1/(2\sqrt d)),\qquad
 \kappa_d=\bar\kappa(1-1/(2\sqrt d)).
\]
Take \(\bar\eta<1\) and \(\bar\kappa>0\) sufficiently small that \(\bar\kappa<\epsilon\bar\eta/4\); both can be fixed as absolute constants. Set \(\kappa_-=\bar\kappa/2\). Fix \(\xi>0\) sufficiently small, e.g. \(\min(1/32,\kappa_-/4)\), and use the small-level estimate with \(\delta=\xi/4\). In the range of \eqref{ten:spin:eq:5}, \(D_\lambda\le n^k\), so \eqref{ten:spin:eq:3} holds with \(J_n(\lambda)\le 1\) for sufficiently large fixed \(r\). For any fixed finite range of sizes, Lemma~\ref{lem:finitegap} and
\(J_n(\lambda)\le D_\lambda^2\|B_n\|_{V_\lambda}^{2r}\)
allow one \(r\) to make every nontrivial weighted moment at most one. The trivial moment is one.

We prove \eqref{ten:spin:eq:3} by induction on \(d\), using a split into nearly equal numbers of coordinates. We only need to perform the induction step for sufficiently large \(d\), with all largeness bounds independent of \(r\); the same \(r\), taken as a positive integer, can then be used throughout. In view of the previous bounds we assume
\begin{equation}
 \lambda_1^t\le n/2,\qquad k=n-\lambda_1>n^{1-\delta},
 \qquad \log D_\lambda\ \ge c n^{1-\delta}.
 \label{ten:spin:eq:6}
\end{equation}
(The trivial case is already covered.) Put \(s=2^{\lfloor d/2\rfloor}\), \(m=2^{\lceil d/2\rceil}\), so positions form \(m\) columns of size \(s\), and \(s\) rows of size \(m\). Take the columns and rows corresponding to the consecutive parts of the sweep. Let \(B_{\mathrm{col}}\) be the product of the size \(s\) column sweeps and \(B_{\mathrm{row}}\) the product of the size \(m\) row sweeps, so
\(B_n=B_{\mathrm{row}}B_{\mathrm{col}}\).
Put \(U=B_{\mathrm{row}}^*B_{\mathrm{row}}\) and
\(V=B_{\mathrm{col}}B_{\mathrm{col}}^*\). The matrices \(B_n^*B_n\) and
\(U^{1/2}VU^{1/2}\) have the same nonzero spectrum: they are \(T^*T\) and \(TT^*\) for \(T=U^{1/2}B_{\mathrm{col}}\). Applying the Araki--Lieb--Thirring trace inequality
\(\operatorname{Tr}(U^{1/2} V U^{1/2})^r\le\operatorname{Tr}(U^r V^r)\) for positive matrices and \(r\ge1\) gives
\begin{equation}
 \operatorname{Tr}_\lambda (B_n^* B_n)^r
 \ \le\ \operatorname{Tr}_\lambda(XY),
 \label{ten:spin:eq:7}
\end{equation}
where \(X=V^r=(B_{\mathrm{col}}B_{\mathrm{col}}^*)^r\) and
\(Y=U^r=(B_{\mathrm{row}}^*B_{\mathrm{row}})^r\) are positive group algebra operators in the columns and rows respectively. They are products over their blocks, with single-block factors \((B_sB_s^*)^r\) in columns and \((B_m^*B_m)^r\) in rows. The two orientations of a child's positive square have the same trace, so both use the child moment in \eqref{ten:spin:eq:3}.

Write $\eta'=\eta_{\lceil d/2\rceil}$ and
$\kappa'=\kappa_{\lceil d/2\rceil}$, the larger coefficients of the
two child sizes.

\subsection{Signed spins and holes}
We can first assume that both \(\alpha,\beta\) have length \(\le q=\lceil n^\epsilon\rceil\). Indeed, truncate each after \(q\) parts and move the remaining boxes into the number \(l\), say increasing it to \(L\). By branching and \eqref{ten:spin:eq:1} the new pair of diagrams works with some diagram on \(L\) positions. If \(F'\) is the new entropy, \(F-F'\ge (L-l)\log q\) since removed parts have size at most \((u+v)/q\). And \(g_n(L)-g_n(l)\le (\kappa_d\log n+1)(L-l)\). Thus the bound with the new data implies the desired one, for large \(d\).

For data with these bounded lengths, our next objective is
\[
 \log J_n(\lambda)
 \le\eta' F(\alpha,\beta)+\kappa'l\log n-l\log(n/l)
                       +O(l+n^{1-\xi}).
\]
The local stabilizer traces will retain the child entropies; the
angle estimate will convert them to the prescribed global entropy.
Counting the labeled placements will produce the negative remainder
term. The increments from $\eta',\kappa'$ to the parent coefficients
will then absorb the displayed error.

Use the signed space $V=\mathbb C^q\oplus\mathbb C^q$ and its
commuting group $U(q)\times U(q)$ from Section~\ref{sec:static}.
For a local spin type $(a,b)$ its permutation representation is
\begin{equation}
 \operatorname{Ind}_{S_{|a|}\times S_{|b|}}^{S_{|a|+|b|}}
        (V_a\otimes V_{b^t})\otimes\mathcal U_a\otimes\mathcal U_b,
 \label{ten:spin:eq:8}
\end{equation}
by \eqref{eq:signed-sector-decomposition}. Its compact carrier has
dimension at most $(n+1)^{q(q-1)}$. Tensor positions may be regrouped
by the signed permutation unitaries; parity-preserving spin matrices
regroup in the usual way.

Allow \(l\) holes in addition, specially marked by distinct labels, one assignment \(z\) specifying distinct sites for these holes. The space is the orthogonal sum over \(z\) of spin spaces on the remaining sites, with holes permuting along with positions (holes count as plus in the sign convention). Project to global spin type \((\alpha,\beta)\) under the simultaneous \(U(q)\times U(q)\); call this projector \(P\), with its diagonal blocks at fixed assignments also denoted \(P\). Its space as a permutation representation contains \(V_\lambda\) with multiplicity at least \(D_\gamma\), by \eqref{ten:spin:eq:8}, \eqref{ten:spin:eq:1}, and the regular permutation space for ordering the \(l\) holes. Since \eqref{ten:spin:eq:7} involves a trace of a product of positive elements, each irreducible contributes nonnegatively, so
\begin{equation}
 D_\gamma \operatorname{Tr}_\lambda(XY)
 \le \operatorname{Tr}(PXY)
 =\sum_z \operatorname{Tr}(P X_z Y_z).
 \label{ten:spin:eq:9}
\end{equation}
Here \(X_z,Y_z\) are the diagonal blocks at \(z\), both positive. This diagonal-block reduction holds because \(X\) preserves the column of every hole label, \(Y\) the row, and hence in the trace the intermediate location must equal the specified location for each label. The projectors for the global spin group leave assignments fixed. At fixed \(z\), \(X_z\) factors in the column spin blocks on remaining sites, and \(Y_z\) similarly in the row spin blocks. Write \(h_j\) for the number of holes in column \(j\), and \(b_i\) for the number in row \(i\).

\subsection{Local stabilizer trace bounds}
We give local trace bounds for \eqref{ten:spin:eq:9}. Use local column spin-unitary type projectors \(P_H\), one type specified in each column, and local row type projectors \(P_K\). For these choices let \(F_H,F_K\) be the sums of the entropies \eqref{ten:spin:eq:2} of their respective types, using each block's own total spin-site count. All these projectors commute with \(P\); \(X_z\) commutes with \(P,P_H\) and \(Y_z\) with \(P,P_K\).
Then for each fixed set of column types,
\begin{equation}
 \operatorname{Tr}(X_z P_H)
 \le
 \left(\prod_j(s)_{h_j}^{-1}\right)
 \exp\left(-(1-\eta')F_H+\sum_j g_s(h_j)+O(n^{1-\xi})\right),
 \label{ten:spin:eq:10}
\end{equation}
where \((s)_h=s(s-1)\cdots(s-h+1)\). There is the analogous row bound with \(m,b_i\).

Apply Lemma~\ref{s:local-stabilizer-trace} to each column with
$t=s$, $h=h_j$, $N=n$, $G=g_s(h_j)$, and its specified local type $(a,b)$.
For every constituent $\rho$ the exact restriction step is
\begin{equation}
 \operatorname{Tr}_\rho(E_{S_{s-h}}x)
 \le\frac{1}{(s)_hD_\rho}\sum_{\mu\supseteq\rho}J_s(\mu).
 \label{ten:spin:eq:11}
\end{equation}
Every extension $\mu$ has an occurrence of $a,b$ together with some
$h$-box diagram, as proved in that lemma. Thus the every-occurrence
induction hypothesis applies, with the larger child coefficients
$\eta',\kappa'$. Lemma~\ref{lem:signed-hook-carriers} gives
\begin{equation}
 \log D_\rho\ge F(a,b)-O(q^2\log(n+1)).
 \label{ten:spin:eq:12}
\end{equation}
The local trace lemma has already included the multiplicities and the
compact carrier; these must not be inserted a second time. Its remaining
factor is the partition count $\mathsf p(s)$.
Multiplying over lines, the logarithms of all carrier factors, partition
counts, and choices of local types total
\[
 O\bigl(n^{1/2+2\epsilon}\log(n+1)
                +n^{3/4}\log(n+1)\bigr)=O(n^{1-\xi}).
\]
The same calculation applies to rows. This proves
\eqref{ten:spin:eq:10}, including full columns of holes and empty
spin types. Every constant is independent of the moment $r$.

\subsection{An angle estimate from support inverses}
Next we need the following angle bound on the type projectors on the spin space at fixed \(z\), writing \(F=F(\alpha,\beta)\):
\begin{equation}
 \|P_H P_K P\|^2
 \le \min\{1,\ \exp(F_H+F_K-F+ l+O(n^{1-\xi}))\}.
 \label{ten:spin:eq:13}
\end{equation}
To apply Theorem~\ref{thm:one-sided-type-angle}, take its occupied board to be the complement of the fixed hole assignment $z$. Its global type is $(\alpha,\beta)$, its column types are the specified $P_H$, and its row types are the specified $P_K$. The local carrier product is at most
\[
 (n+1)^{(s+m)q(q-1)}=\exp(O(n^{1/2+2\epsilon}\log(n+1))).
\]
This is absorbed by $O(n^{1-\xi})$ because $\epsilon=1/64$ and $\xi\le1/32$. The theorem's missing-cell factor is at most $e^l$, with no assumption on the arrangement of the holes. Its entropy terms are exactly $F_H,F_K,F$. This proves~\eqref{ten:spin:eq:13}.

\subsection{Combining the two positive traces}
Using positivity, expansion into pairs of types in \eqref{ten:spin:eq:9} bounds each term by
\[
 \|P_H P_KP\|^2\,\operatorname{Tr}(X_z P_H)\operatorname{Tr}(Y_z P_K).
\]
Indeed spectral decompositions of the positive restrictions (also restricting to the global projector) give the angle squared as an upper bound for squared scalar products. Using \eqref{ten:spin:eq:10}, its row version, and \eqref{ten:spin:eq:13} gains the exponent \(-(1-\eta')F\), since for \(w=F_H+F_K\)
\[
 -(1-\eta') w+\min(0,w-F)\le -(1-\eta')F .
\]
Summing over type choices within our log errors, we obtain
\begin{equation}
 \begin{aligned}
 \operatorname{Tr}(P X_z Y_z)
 &\le \prod_j(s)_{h_j}^{-1}\prod_i(m)_{b_i}^{-1}\\
 &\quad\times\exp\left(-(1-\eta')F+\sum_j g_s(h_j)+\sum_i g_m(b_i)
                     +O(l+n^{1-\xi})\right).
 \end{aligned}
 \label{ten:spin:eq:14}
\end{equation}

\subsection{Entropy of the hole assignments}
We detail the count over hole assignments; it is important to keep the entropy terms from \(g_s,g_m\). First the falling factorials give a bound \(n^{-l} e^{O(l)}\) for the prefactor (one can use \((a)_p\ge a^p e^{-O(p)}\) for \(p\le a\)). Also
\begin{equation}
 \sum_j g_s(h_j)+\sum_i g_m(b_i)
 \le \kappa' l\log n
   -\sum_j h_j\log(s/h_j)-\sum_i b_i\log(m/b_i)
   +O(n^{1-\xi}).
 \label{ten:spin:eq:15}
\end{equation}
In fact taking a maximum with zero in a column expression can raise it by at most \(O(s^{1-\kappa_-}\log n)\) since it is only needed at \(h_j\le s^{1-\kappa_-}\); use the similar row error. Finally
\begin{equation}
 \sum_z n^{-l}\exp\left(-\sum_j h_j\log(s/h_j)-\sum_i b_i\log(m/b_i)\right)
 \le \exp\left(-2l\log(n/l)+O(n^{1-\xi})\right).
 \label{ten:spin:eq:16}
\end{equation}
Drop distinctness, so normalized counting uses independent column and row assignments with replacement to cells. For \(l>0\) the column negative sum exponent is \(-l\log(n/l)+l D(h/l\ \|\ \mathrm{unif}_m)\) with \(D\) the relative entropy of proportions (\(h=(h_j)\)). Each possible histogram has probability at most the exponential of minus the divergence term \(lD\), by the multinomial formula. The row count is analogous and independent, and the log numbers of histograms fit the error. More explicitly the normalized sum is at most
\[
 e^{-2l\log(n/l)}
 \binom{l+m-1}{m-1}\binom{l+s-1}{s-1}.
\]
The logarithm of the two histogram counts is
$O((s+m)\log(n+1))$, within the stated error. For $l=0$ the sum and
the histogram product are one. All formulas omit zero-count terms.

Using \(D_\lambda\le \binom n l e^F D_\gamma\) by \eqref{ten:spin:eq:1}, we conclude from \eqref{ten:spin:eq:7}, \eqref{ten:spin:eq:9}, \eqref{ten:spin:eq:14}--\eqref{ten:spin:eq:16} that
\begin{equation}
 \log J_n(\lambda)
 \le \eta' F+\kappa'l\log n-l\log(n/l)+O(l+n^{1-\xi}).
 \label{ten:spin:eq:17}
\end{equation}
It remains to absorb the error in \eqref{ten:spin:eq:17}. Put
$\Delta\eta=\eta_d-\eta'$ and $\Delta\kappa=\kappa_d-\kappa'$.
Both are bounded below by a positive constant times $d^{-1/2}$ for
large $d$. By \eqref{ten:spin:eq:6} and the same dimension upper
bound, now using $D_\gamma\le l!$, we have
$F+l\log n\ge c n^{1-\delta}$. The available increment can be
divided into two parts:
\[
 \tfrac12\Delta\kappa\,l\log n\ge c l\sqrt d,
 \qquad
 \Delta\eta F+\tfrac12\Delta\kappa\,l\log n
 \ge c\frac{n^{1-\delta}}{\sqrt d}.
\]
The first dominates $O(l)$, and the second dominates
$O(n^{1-\xi})$ because $\delta<\xi$. This proves the desired bound:
replacing $\kappa_d l\log n-l\log(n/l)$ by its maximum with zero
only increases the right-hand side.

The order of choices is now explicit. First choose a threshold $d_0$
large enough for every large-size comparison above; all its constants
are independent of $r$. Next choose one integer $r$ exceeding the sparse
threshold and all finitely many thresholds from
Lemma~\ref{lem:finitegap} for $d\le d_0$. Enlarging $r$ decreases
all moments. Together with the truncation argument, this completes the
induction proving \eqref{ten:spin:eq:3}.

\subsection{Full-deck amplification}
The every-occurrence estimate is now proved. To turn it into a bound
for the whole permutation law, we choose a zero-hole occurrence whose
entropy is controlled by the irreducible dimension.

\begin{lemma}\label{lem:spin-zero-hole}
There is an absolute $C_0$ such that, for all sufficiently large $n$,
every $\lambda\vdash n$ with at most $n/2$ rows has an occurrence
\eqref{ten:spin:eq:1} with $l=0$ and
\begin{equation}
 F(\alpha,\beta)\le C_0\log D_\lambda.
 \label{ten:spin:eq:4}
\end{equation}
\end{lemma}
\begin{proof}
Taking any number of top rows as $\alpha$ and the remaining diagram
as $\beta^t$ gives an occurrence by the Littlewood--Richardson rule.
First split at the Durfee square size $a$. Write $x_i$, $1\le i\le a$,
for the top row lengths and $y_j$, $1\le j\le a$, for the column
lengths below those rows. The hook product is at most
\[
 \left(\prod_i x_i!\prod_j y_j!\right)
 \left(\prod_i\binom{x_i+a}{a}
             \prod_j\binom{y_j+a}{a}\right)
 \prod_j(1+y_j/a)^a.
\]
Indeed, ignoring cells below the top $a$ rows, a top-row hook is
at most its horizontal length plus $a$. Inside the square the
additional bottom column length costs the factor $1+y_j/a$.
The bottom hooks satisfy the transposed bound.
The logarithm of the first binomial product is at most
\[
 a\sum_i\log\bigl(e(1+x_i/a)\bigr)
 \le a^2\log\bigl(e(1+n/a^2)\bigr)=O(n),
\]
using $a^2\le n$. The other binomial product has the same bound, and
$\sum_j a\log(1+y_j/a)\le\sum_jy_j\le n$.
The hook formula and factorial estimates therefore give
$F(\alpha,\beta)\le\log D_\lambda+O(n)$.

Put $k=n-\lambda_1$. If $k\ge n/4$, Lemma~\ref{lem:spin-dimension-range}
absorbs this $O(n)$ term. If $1\le k<n/4$, instead apply the Durfee
split to the tail $\bar\lambda\vdash k$ and include the first row
in $\alpha$. The resulting entropy is at most
\[
 k\log(n/k)+k+\log D_{\bar\lambda}+O(k).
\]
Here the first two terms bound the entropy of separating the first
row from the tail. Tableaux of the tail extend to $\lambda$, so
$D_{\bar\lambda}\le D_\lambda$, and the first part of
Lemma~\ref{lem:spin-dimension-range} bounds the other terms by
$O(\log D_\lambda)$. Finally, $k=0$ gives entropy zero.
\end{proof}

\begin{proof}[Proof of Corollary~\ref{cor:spin-full-deck}]
Choose \(\bar\eta\) so that \(C_0\bar\eta<1/2\) in \eqref{ten:spin:eq:4}. Taking \(l=0\) with those choices in \eqref{ten:spin:eq:3}, every nontrivial representation not already annihilated satisfies
\[
 \|B_n\|_{V_\lambda}^{2r}\le D_\lambda^{-1/2}
\]
for large \(n\). Lemma~\ref{lem:degrees} now provides the required Fourier sum.
After $t=6r$ forward sweeps, nonnormality causes no problem:
$\|B_n^t\|_{\mathrm{HS}}^2\le D_\lambda\|B_n\|^{2t}$ in each
irreducible. Consequently the squared $L^2$ distance of the density
from uniform is at most
\[
 \sum_{\lambda\ne(n)}D_\lambda^2\|B_n(\lambda)\|^{12r}
 \le\sum_{\lambda\ne(n),(1^n)}D_\lambda^{-1}=o(1).
\]
The sign representation and all shapes with too many rows contribute
zero by Lemma~\ref{lem:annihilation}. Translation invariance makes the
bound uniform over initial orders. Since one sweep is $d$ physical
shuffles, $6rd$ shuffles suffice at total variation $1/4$ for all
sufficiently large $d$. Together with the support obstruction, this
proves the stated $\Theta(d)$ mixing consequence.
\end{proof}

\paragraph{Local access to a random permutation.}
The full-deck estimate also gives a short decision-tree description of
an approximately uniform permutation. A decision forest computes the
coordinates of its output by separate adaptive decision trees that
query a shared random input. This is the permutation-sampling model
studied by Alekseev, G\"o\"os, Myasnikov, Riazanov, and
Sokolov~\cite[Section~1]{AlekseevEtAl2025}.

\begin{corollary}[Permutation sampling with logarithmic query depth]
\label{cor:decision-forest}
For $n=2^d$ there is a decision forest whose output is always a
permutation of $[n]$, whose input consists of independent fair bits,
and whose coordinate trees have depth $O(\log n)$, such that the
output law has total variation distance $o(1)$ from uniform as
$d\to\infty$. The same conclusion holds with independent uniform
$[n]$-valued input cells.
\end{corollary}
\begin{proof}
Use the fixed number $L$ of sweeps in
Corollary~\ref{cor:spin-full-deck}. Assign one input bit to every
switch in the resulting $Ld$ layers, with all coordinate trees
using the same assignment. To compute the final position of card
$i$, follow its path through the layers. At each layer the current
position determines which switch bit to query, and that bit determines
the next position. Thus each tree makes $Ld$ adaptive queries. For
every input the switches compose to a permutation; jointly, the
outputs have exactly the $L$-sweep law. The claimed convergence follows
from Corollary~\ref{cor:spin-full-deck}. For $[n]$-valued input cells,
use the parity of one independent cell for each switch; it is a fair
bit because $n$ is even.
\end{proof}

For dyadic sizes, this improves the $O((\log n)^2)$ shared-cell
upper bound recalled in~\cite[Section~1]{AlekseevEtAl2025} to
$O(\log n)$. The same depth order holds in the fair-bit model.
The depth order is optimal for fair-bit queries at any fixed error
below one: a binary tree of depth $h$ has at most $2^h$ output values,
so its marginal is at total variation distance at least $1-2^h/n$
from uniform on $[n]$. This argument concerns bit queries; it does
not give the same lower bound for $[n]$-valued cell queries.
